% Source: Wald GR, physical PDF p. 205
%         (printed p. 192).
% Coverage: A set and the boundary of its chronological future.
% Figures/tables/footnotes: Figure 8.3; no tables or footnotes.
% Status: source-reviewed and compile-clean.
% Uncertainties: none.

\begingroup
\renewcommand{\thefigure}{8.3}
\begin{figure}[H]
  \centering
  \begin{tikzpicture}[x=.88cm,y=.88cm,>=Latex,line cap=round,font=\small]
    \draw[thick] (-3.15,2.00) -- (-1.42,.20)
      .. controls (-.54,.54) and (.34,.26) .. (1.18,.38) -- (2.84,1.88);
    \draw[->] (-1.16,1.18) -- (-2.12,1.02);
    \node at (-.80,1.30) {\(\partial I^+(S)\)};
    \draw[->] (.54,-.10) -- (.43,.35) node[below=.15cm] {\(S\)};
    \fill (1.89,1.35) circle (1.6pt) node[right] {\(q\)};
    \fill (1.91,2.34) circle (1.6pt) node[right] {\(p\)};
    \draw (1.89,1.35) .. controls (1.65,1.11) and (1.54,1.43) .. (1.63,1.58)
      .. controls (1.77,1.79) and (2.10,1.63) .. (2.18,1.39)
      .. controls (2.25,1.19) and (2.08,1.05) .. (1.89,1.35);
    \draw[->] (1.48,1.83) -- (1.67,1.53);
    \node at (1.35,1.95) {\(O\)};
  \end{tikzpicture}
  \caption{A spacetime diagram showing a set \(S\) and the boundary of its future,
  \(\partial I^+(S)\) (see theorem~\ref{thm:8.1.3}).}
  \label{fig:8.3}
\end{figure}
\endgroup
